%& -shell-escape
% !TeX root = thesis-abstract.tex
% !TeX encoding = UTF-8
% !TeX program = XeLaTeX

\documentclass{.local/lib/classes/ndvr-super-article-standalone}

\begin{document}

    В кинематографе существует три различных понятия. 
    Кадр или фотографический кадр — статическое изображение. 
    В мультипликации для его обозначения используют термин «кадрик». 
    Сцена — множество кадров, связанных единством места и времени. 
    Съёмка или монтажный план — множество кадров, 
    связанных единством процесса съёмки. 
    Сцена может включать несколько съёмок\footnote{
        Fabro М.D., Böszörmenyi L. State-of-the-art 
        and future challenges in video scene detection: a survey
        In Multimedia Systems, Vol. 19, № 5; (2013), 
        pp. 427-454, doi:10.1007/s00530-013-0306-4 
        by Manfred Del Fabro, Laszlo Böszörmenyi, 
        Multimedia Systems Volume 19, Issue 5 , pp 427-454 
    }.

    \begin{definition}
    Съёмка или кинематографический кадр~---~совокупность 
    множества фотографических кадров $\videoframe$
    внутри временной области $\videoline$, 
    кадры которой $\videoframe_{ \videoshot, i}$
    значительно отличаются от кадров соседних областей.
    \[
        \left\lbrace
            \begin{array}{rll}
                \videoshot & = &  \{
                        \videoframe_{ \videoshot, i}
                            ; \videoframediff(\videoframe_{ \videoshot, i},
                                \videoframe_{\videoshot, j})
                                    < { \varepsilon},
                                \videoframe_{\videoshot, i},
                                \videoframe_{\videoshot, j}
                                \in \videoline
                    \}\\
                \videoframediff & = &  \text{функция разности кадров}.

            \end{array}
        \right.
    \]
    \end{definition}
    \begin{definition}
        Аномалия — скачкообразное изменение свойств 
        наблюдаемого ряда в заранее неизвестный момент времени.
        Иногда аномалию называют <<разладкой>>.
    \end{definition}
    \noindent Будем называть моментами событий~---~моменты времени,
    в~которых наблюдались аномалии. 

    \begin{statement}[О событиях]
        Моменты событий совпадают с моментами смены съёмок. 
    \end{statement}

\end{document}
 
